\section{Related Work}

XXX:  This is pretty messy.  It's pretty much just random thoughts now.

We will focus on overflow bugs.

It is well known that C code tends to have security holes due to buffer
overrun bugs and integer overflow bugs.  Many approaches have been applied
to preventing, or at least detecting, these bugs.  While several have shown
a substantial improvement, often with a minimal cost of efficiency, so far
none of them have provided a solution that solves the problem in all cases.
In fact, there is little hope that any solution can completely protect
existing C code from the bugs it contains; ultimately, if the code has a
bug, the best we can hope for (without modifying the language) is to limit
the damage caused by it.  We would prefer to write code that is free of
bugs in the first place, but experience has proven that it is simply
infeasible to do this .

The solution that Pit provides does come at a cost; however, the cost is
optional and can be mitigated substantially with compiler optimizing
techniques.

, the tradeoff between security and efficiency still exists.  We find,
however, that most of the time, this tradeoff is not as bad as it seems.  A
tremendous amount of time is spent in a small fraction of a program's code.
(XXX: find work to support this)  Therefore, we can easily reduce the
speed of most of our code by a factor of, for example, five or ten with no
noticeable impact, as long as the critical sections, that is sections where
most of the time is spent, can still be written efficiently.  If we can
easily guarantee that the code outside of the critical sections has no
buffer overrun bugs or integer overflow bus, we can concentrate our efforts
on making sure the critical sections are both secure and efficient.

\begin{itemize}
\item
Perl.  Perl code never has buffer overflows (explain why).  Also code
is easier to write because of managed arrays and hashes.
\item
Articles about buffer and integer overflow vulnerabilities
\item
Safe Systems Programming Languages (wpeii.pdf) is worth referencing
here.  [[Safe Systems Programming Languages (wpeii.pdf) is
interesting, and I'm reading it more closely.  It seems to relate
really well.  I like their goals for a programming language for
systems; this describes a good distinction between languages like C
and Pit and higher ones like Java.  It also introduced me to some
interesting ways of protecting existing C code and a couple languages
I hadn't heard of yet.  Start reading again at 4.2.]]
\end{itemize}
